\subsection{CASIMIR 元}

根据第 1 节第 3 号命题 3，g 上存在负的、对称的、分离的且在 Ad(G) 下不变的双线性型（例如，若 G 是半单的，可以取 g 的 Killing 型）。令 F 为这样的一个型。回忆（第 I 章，第 3 节，第 7 号），与 F 关联的 Casimir 元是包络代数 U(g) 的中心元素 \(\Gamma\)，满足：对于 g 的任意基 \((e_{i})\)，若 \(\mathrm{F}(e_{i}, e_{j}) = -\delta_{ij}\)，则有 \(\Gamma = -\sum e_{i}^{2}\)。

在本章余下部分中，我们把 G 的 Casimir 元称为由 g 上负的、分离的、不变的对称双线性型按上述方式得到的 \(\mathrm{U}(\mathfrak{g})\) 的元素。若 \(\Gamma\) 是 G 的 Casimir 元，且 \(\tau: G \to \mathbf{GL}(V)\) 是 G 的不可约表示，则 V 的自同态 \(\Gamma_{V}\) 是一个伸缩变换（《代数》，第 VIII 章，第 3 节，第 2 号，定理 1），其比率记为 \(\tilde{\Gamma}(\tau)\)。

命题 4. 令 \(\Gamma\) 为 G 的 Casimir 元。

\begin{enumerate}
\def\labelenumi{\alph{enumi})}
\item
  若 \(\tau\) 是 \(\mathbf{G}\) 的不可约表示，则 \(\tilde{\Gamma}(\tau)\) 为正实数。若 \(\tau\) 不是平凡表示，则 \(\tilde{\Gamma}(\tau) > 0\)。
\item
  存在唯一的二次型 \(\mathbf{Q}_{\Gamma}\)，定义在 \(\mathrm{X(T)}\otimes \mathbf{R}\) 上，使得对于每个不可约表示 \(\tau\)，它属于 \(\mathbf{G}\)，
\end{enumerate}

\[
\tilde {\Gamma} (\tau) = \mathrm{Q} _ {\Gamma} (\lambda + \rho) - \mathrm{Q} _ {\Gamma} (\rho),
\]

其中 \(\lambda\) 是 \(\tau\) 的最高权。型 \(Q_{\Gamma}\) 是正的、分离的且在 W 下不变。

令 \(\mathrm{F}\) 为 \(\mathfrak{g}\) 上定义 \(\Gamma\) 的分离负对称双线性型。令 \(\tau : \mathrm{G} \to \mathbf{GL}(\mathrm{V})\) 为 \(\mathrm{G}\) 的不可约表示；令 \(\langle , \rangle\) 为 \(\mathrm{V}\) 上在 \(\mathrm{G}\) 下不变的 Hilbert 内积（第 1 节，第 1 号），并令 \((e_i)\) 为 \(\mathfrak{g}\) 的满足 \(\mathrm{F}(e_i, e_j) = -\delta_{ij}\) 的基。则对于每个元素 \(v\)，它属于 \(\mathrm{V}\) 且不在 \(\mathrm{G}\) 下不变，有

\[
\begin{array}{c} \tilde {\Gamma} (\tau) \langle v, v \rangle = \langle v, \Gamma_ {\mathrm{V}} (v) \rangle = - \sum_ {i} \langle v, (e _ {i}) _ {\mathrm{V}} ^ {2} v \rangle = \sum_ {i} \langle v, (e _ {i}) _ {\mathrm{V}} ^ {*} (e _ {i}) _ {\mathrm{V}} v \rangle \\ = \sum_ {i} \langle (e _ {i}) _ {\mathrm{V}} v, (e _ {i}) _ {\mathrm{V}} v \rangle > 0, \end{array}
\]

故 \(a\) 成立。

令 B 为 \(t_{C}^{*}\) 上的逆型，它是限制到 \(t_{C}\) 的、由 F 经扩张标量从 \(g_{C}\) 上诱导的双线性型的逆型。根据第 VIII 章第 6 节第 4 号命题 7 的推论，有 \(\tilde{\Gamma}(\tau) = \mathrm{B}(\delta(\lambda), \delta(\lambda + 2\rho))\)。将 \(\delta : \mathrm{X(T)} \to t_{\mathbf{C}}^{*}\) 延拓为从 \(\mathrm{X(T)} \otimes \mathbf{R}\) 到 \(t_{C}^{*}\) 的 R-线性映射，并令 \(Q_{\Gamma}\) 为二次型 \(x \mapsto \mathrm{B}(\delta(x), \delta(x))\)，定义在 \(\mathrm{X(T)} \otimes \mathbf{R}\) 上；它是分离的且在 W 下不变，并且

\[
\tilde {\Gamma} (\tau) = \mathrm{B} (\delta (\lambda + \rho), \delta (\lambda + \rho)) - \mathrm{B} (\delta (\rho), \delta (\rho)) = \mathrm{Q} _ {\Gamma} (\lambda + \rho) - \mathrm{Q} _ {\Gamma} (\rho).
\]

我们证明型 \(Q_{\Gamma}\) 是正的。事实上，若 \(x \in \mathrm{X}(\mathrm{T}) \otimes \mathbf{Q}\)，则元素 \(\delta(x)\) 属于 \(t_{C}^{*}\)，它在 t 上取纯虚值，因而在 it 上取实值；注意到对于 \(y \in it\)，有 \(\mathrm{F}(y, y) \geq 0\)，即可得出结论。

还需证明 b) 中的唯一性断言。令 Q 为 \(X(T) \otimes \mathbf{R}\) 上满足所需条件的二次型，令 \(\Phi\)（分别为 \(\Phi_{\Gamma}\)）为与 Q（分别与 \(Q_{\Gamma}\)）关联的双线性型。对于 \(\lambda, \mu \in X_{++}\)，有

\[
\begin{array}{l} \Phi (\lambda , \mu) = (Q (\lambda + \mu + \rho) - Q (\rho)) - (Q (\lambda + \rho) - Q (\rho)) - (Q (\mu + \rho) - Q (\rho)) \\ \qquad = \Phi_ {\Gamma} (\lambda , \mu). \end{array}
\]

由于 \(\mathrm{X(T)_{++}}\) 生成 R-向量空间 \(\mathrm{X(T)}\otimes\mathbf{R}\)，有 \(\Phi=\Phi_{\Gamma}\)，所以 \(Q=Q_{\Gamma}\)。

注. 令 \(x \in g\)。存在严格正实数 A，使得对于每个不可约表示 \(\tau : G \to \mathbf{GL}(V)\) 以及 V 上每个在 G 下不变的 Hilbert 结构，

\[
\| \mathrm{L} (\tau) (x) \| ^ {2} \leq \mathrm{A}. \tilde {\Gamma} (\tau).
\]

事实上，沿用前面证明中的记号，可以选取 g 的基 \((e_{i})\)，使得 \(x = ae_{1}, a \in R\)。于是对于 \(v \in V\)，有

\[
\langle x _ {\mathrm{V}} v, x _ {\mathrm{V}} v \rangle = | a | ^ {2} \langle e _ {1} v, e _ {1} v \rangle \leq | a | ^ {2} \tilde {\Gamma} (\tau) \langle v, v \rangle .
\]

